\chapter{Optical and Quasiparticle Properties}
\label{ch:optics-gw}

Four workflows evaluate the response of a converged ground state: the
dielectric function, the same response reduced to sheet quantities for a 2D
material, the second-order (nonlinear) response, and the $G_0W_0$ quasiparticle
correction.

\begin{calwarn}
The first two and $G_0W_0$ \textbf{inherit a baseline} rather than converging
their own ground state --- a completed GPAW
\menu{Simulation}\menusep\menu{Single-point Calculation\dots}, which writes
\file{single\_point.gpw}. They will not open without one.

This is not merely a saving. Re-converging inside the job would give a spectrum
from a different SCF solution than the one you validated --- different
smearing, a different $k$-grid, possibly a different magnetic state --- with
nothing in the output to say so. Because the ground state is inherited whole,
these wizards show no Calculator Settings stage: the cutoff, functional and
$k$-grid all come back from the baseline, and offering to set them again would
present knobs that cannot affect the run.

Nonlinear Optics (\S\ref{sec:nlopt}) is the exception, and for a reason worth
knowing: \texttt{gpaw.nlopt} requires point-group symmetry to be off and its
band sums need a \emph{converged} empty manifold, neither of which a general
baseline has. It therefore converges its own ground state and keeps a
Calculator Settings stage.
\end{calwarn}

\section{Optical properties}
\label{sec:optics}

\menu{Simulation}\menusep\menu{Optics\dots} evaluates the frequency-dependent
dielectric function $\varepsilon(\omega)$ with GPAW's linear-response module
and derives the usual optical observables from it.

\subsection{Settings}

\begin{description}
  \item[\ui{Baseline SCF (.gpw)}] The ground state to inherit. The note beneath
    it spells out what is being inherited --- functional, cutoff, $k$-grid and
    environment --- because with no Calculator Settings stage this is the only
    place those values are visible, and they are what decides whether the
    spectrum is trustworthy.
  \item[\ui{Broadening $\eta$}] Lorentzian broadening, default 0.1\,eV.
    Smaller values resolve sharper features but need a denser $k$-mesh in the
    \emph{baseline} to be meaningful.
  \item[\ui{Energy window min/max}] The photon-energy range, default
    0--20\,eV.
  \item[\ui{Tetrahedron integration}] See below.
  \item[\ui{Number of points}] Samples on the energy grid.
  \item[\ui{Polarization directions}] $\varepsilon_{xx}$, $\varepsilon_{yy}$,
    $\varepsilon_{zz}$. For an isotropic crystal the three coincide.
\end{description}

\subsection{Brillouin-zone integration}

By default the response is summed over $k$-points, giving every interband
transition a Lorentzian of width $\eta$. Anything narrower than $\eta$ --- a
van Hove singularity, a 2D absorption edge --- is smeared into the broadening
rather than resolved.

\ui{Tetrahedron integration} interpolates the bands linearly within each
tetrahedron and integrates analytically, resolving those features on a mesh
where point integration is still noisy.

\begin{calwarn}
Tetrahedron integration requires the \emph{baseline's} $k$-grid to contain
every vertex of the irreducible Brillouin zone --- a grid from
\opt{gpaw.bztools.find\_high\_symmetry\_monkhorst\_pack()}. An ordinary
Monkhorst--Pack grid usually does not qualify, and since the grid is inherited,
it has to be right when the \emph{baseline} is run; the optics job cannot fix
it afterwards. When the grid is unsuitable the run stops and says so, naming
the remedy, rather than quietly falling back to point integration and returning
a spectrum computed by a different method than the one requested. Which
integrator ran is recorded in \file{optics.json} alongside the spectra.
\end{calwarn}

\subsection{The viewer}

The results window plots one quantity at a time against photon energy:
$\varepsilon_1$ and $\varepsilon_2$ together, the absorption coefficient
$\alpha(\omega)$, reflectivity $R(\omega)$, the refractive index $n$ and $k$,
and the energy-loss function $L(\omega)$.

The \ui{X axis} selector switches between \ui{Energy (eV)} and
\ui{Wavelength (nm)} via $\lambda = hc/E$ with $hc = 1239.84197$\,eV\,nm. Every
curve and every tick label follows.

\begin{calnote}
Samples at $E \le 0$ have no finite wavelength and are omitted from the
wavelength view --- dropped from the abscissa and from every curve together, so
nothing is shifted out of alignment. Since GPAW's energy grid starts at zero,
expect the wavelength view to begin one sample in. Very low energies map to
very long wavelengths, so a window that starts near 0\,eV will stretch far into
the infrared; set \ui{Energy window min} above zero if that compresses the part
of the spectrum you care about.
\end{calnote}

\ui{Export CSV\dots} writes one row per sample with every quantity for the
selected direction; \ui{Export Image\dots} renders the current plot at high
resolution.

\section{2D optics}
\label{sec:optics-2d}

\menu{Modules}\menusep\menu{2D Materials}\menusep\menu{2D Optics\dots} runs the
same response and then divides the vacuum back out.

The reason is that a slab supercell's $\varepsilon_{3D}$ is not a property of
the material: double the vacuum padding and it moves, because the sheet's
polarization is diluted over more cell. The wizard adds one question ---
\ui{Vacuum axis}, seeded from the cell but worth confirming, since getting it
wrong rescales every 2D quantity by the wrong length --- and derives three
quantities that do not depend on the padding:

\begin{center}
\small
\begin{tabular}{@{}llp{6.2cm}@{}}
\toprule
Quantity & Definition & Units \\
\midrule
Sheet polarizability & $\alpha_{2D} = \dfrac{L_z}{4\pi}(\varepsilon_{3D} - 1)$
  & \AA{} \\[6pt]
Absorbance & $A(\omega) = \dfrac{\omega L_z}{c}\,\mathrm{Im}[\varepsilon_{3D}]$
  & dimensionless \\[6pt]
Sheet conductivity & $\sigma_{2D} = -i\,\omega\,\alpha_{2D}$
  & $e^2/h$ \\
\bottomrule
\end{tabular}
\end{center}

The $-1$ in $\alpha_{2D}$ removes the vacuum's own contribution, which is what
makes the result a property of the sheet.

The viewer offers these three only when the file actually carries them, and a
2D job opens on the absorbance rather than on $\varepsilon$.

\begin{caltip}
Graphene is the reference case. Its low-energy absorbance is the universal
$A = \pi\alpha \approx 2.29\,\%$, fixed by the fine-structure constant alone,
and its conductivity is the equally universal $\pi/2$ in units of $e^2/h$. If a
graphene run reproduces those, the whole chain --- baseline, response,
vacuum-axis choice and unit conversion --- is working. Reaching the plateau
needs a dense mesh: it is the $k$-sampling near the Dirac point that decides
this, not the cutoff or the band count.
\end{caltip}

\section{Nonlinear optics}
\label{sec:nlopt}

\menu{Electronics}\menusep\menu{Nonlinear Optics\dots} evaluates the
\emph{second-order} response with GPAW's \texttt{gpaw.nlopt} module, within the
independent-particle approximation. Three quantities are offered:

\begin{description}
  \item[\ui{Second-harmonic generation}] $\chi^{(2)}(-2\omega; \omega,
    \omega)$: two photons of energy $\hbar\omega$ in, one of $2\hbar\omega$
    out. Reported in pm/V, and for a sheet additionally as $\chi^{(2)}L$ in
    nm\textsuperscript{2}/V.
  \item[\ui{Shift current}] $\sigma^{(2)}(0; \omega, -\omega)$, in
    A/V\textsuperscript{2}: the DC photocurrent a homogeneous illuminated
    crystal carries with no junction and no built-in field --- the bulk
    photovoltaic effect.
  \item[\ui{Linear susceptibility tensor}] $\chi^{(1)}(\omega)$, the full
    $3\times3$ tensor from the same matrix elements, and the
    $\varepsilon = 1 + \chi^{(1)}$ that follows from it.
\end{description}

\begin{calwarn}
$\chi^{(2)}$ and the shift current are \textbf{odd-rank tensors}: in a
centrosymmetric crystal every component vanishes \emph{identically}, by
symmetry rather than by smallness. GPAW will still return numbers --- the
residue of an incomplete cancellation over a finite $k$-mesh --- and they look
like a spectrum. The generated script therefore tests the cell for an inversion
centre \emph{before} converging anything and warns, and the viewer repeats the
warning above the plot. This is the single most common way the module is
misused.
\end{calwarn}

\subsection{Why this one converges its own ground state}

Unlike every other workflow in this chapter, Nonlinear Optics does \textbf{not}
inherit a Single-Point baseline. Three requirements make an ordinary baseline
unusable:

\begin{itemize}
  \item \texttt{make\_nlodata} asserts that \emph{point-group symmetry is off}.
    The matrix elements it builds are not invariant under the point-group
    folding of the $k$-set, and the assertion arrives with no message --- after
    the SCF has already been paid for.
  \item The sums run over \emph{intermediate} states, so the band set has to be
    large. The wizard's ground state uses \texttt{nbands="nao"}.
  \item Those empty bands must be \emph{converged}. An SCF converges occupied
    states; the unoccupied manifold it leaves behind is noise, and here it is
    summed over. The generated calculator asks for
    \texttt{convergence=\{"bands": -10\}}.
\end{itemize}

All three are imposed by the generator rather than left on the Calculator
Settings stage, because they are requirements of the method and not
preferences. Time reversal is deliberately \emph{kept}: the weaker
\texttt{symmetry="off"} would satisfy the assertion too, at roughly double the
$k$-points, for nothing.

\subsection{Settings}

\begin{description}
  \item[\ui{Tensor components}] Three letters from \texttt{xyz}, several
    separated by commas (\texttt{yyy, xxy}). There is deliberately no
    ``all 27'': most vanish by symmetry, each costs a full sum over bands and
    $k$-points, and 27 spectra of which 24 are numerical noise is not a better
    answer. \menu{Analysis}\menusep\menu{Symmetry} will name the point group
    that decides which are allowed.
  \item[\ui{SHG gauge}] \ui{Length} (GPAW's default) and \ui{Velocity} are
    formally equivalent and numerically are not: the velocity gauge carries
    low-frequency divergences that cancel only for a complete band set, so a
    truncated sum leaves it visibly wrong as $\omega \to 0$. Running both and
    overlaying them is the standard convergence test --- their disagreement
    says more about the band summation than any single number does.
  \item[\ui{Broadening $\eta$}] A second-order spectrum is far more sensitive
    to it than a linear one: the resonances are sharper and the divergences at
    $\omega$ and $2\omega$ are regularized by exactly this number, so a small
    $\eta$ on a coarse mesh produces spikes rather than structure.
  \item[\ui{Scissors shift}] A rigid shift of the empty bands, for the
    semilocal gap error. $\chi^{(2)}$ is more sensitive to the gap than the
    linear response is, because the two-photon resonance sits at \emph{half} of
    it. The applied value is recorded in the results and restated by the
    viewer, so a spectrum never silently carries one.
  \item[\ui{Vacuum axis}] As in 2D optics: a supercell $\chi^{(2)}$ is diluted
    by whatever vacuum was used, and multiplying the thickness back in gives
    the sheet value the 2D literature quotes.
\end{description}

\begin{calnote}
The expensive step is \texttt{make\_nlodata} --- the momentum matrix elements
--- and it runs \emph{once}. Every component and every response reuses the same
data, which is saved as \file{mml.npz} alongside the run. Asking for a second
tensor component costs a band sum, not another ground state.
\end{calnote}

\begin{caltip}
Remember that SHG probes the band structure at $2\omega$ as well as at
$\omega$. A window ending at 6\,eV is sampling transitions up to 12\,eV, and
the band count has to reach them --- which is also why the $\chi^{(1)}$ panel
is worth computing: the $\chi^{(2)}$ features of a semiconductor sit at the
absorption edge \emph{and} at half of it, and a resonance assigned to the wrong
one of those is the standard way to misread an SHG spectrum.
\end{caltip}

\section{GW calculations}
\label{sec:gw}

\menu{Simulation}\menusep\menu{GW Calculations\dots} computes one-shot
$G_0W_0$ quasiparticle corrections,
\[
  E^{\text{qp}}_{n\mathbf{k}} = \varepsilon^{\text{DFT}}_{n\mathbf{k}}
  + Z_{n\mathbf{k}}\left[\Sigma_{n\mathbf{k}}(\varepsilon^{\text{DFT}})
  - V^{xc}_{n\mathbf{k}}\right],
\]
i.e.\ the DFT eigenvalue with its exchange--correlation potential replaced by
the energy-dependent, non-local self-energy. Nothing here is self-consistent,
and the result is only meaningful relative to the baseline it corrects.

\subsection{Engines}

Two engines, each bound to the DFT code that produced its baseline. They are
not interchangeable --- $G_0W_0$ perturbs a specific DFT solution, so the
engine follows from which baseline exists rather than from preference.

\begin{description}
  \item[GPAW] Corrects a \file{.gpw}. The run restarts the baseline, adds a
    generous set of empty bands at fixed density, then applies
    \opt{gpaw.response.g0w0.G0W0}.
  \item[Yambo] Corrects a Quantum ESPRESSO \file{.save} directory. The run
    converts it with \opt{p2y}, generates the $G_0W_0$ input, executes it under
    MPI, and parses the quasiparticle report.
\end{description}

\subsection{Settings}

\begin{description}
  \item[\ui{Screening cutoff}] The convergence parameter that matters. A
    $G_0W_0$ gap is not converged until it stops moving with this.
  \item[\ui{Frequency treatment}] \ui{Plasmon-pole} approximates the frequency
    dependence of the screening with a single pole --- much cheaper, and
    usually good for the gap of a simple semiconductor. \ui{Real axis} performs
    the full integration.
  \item[\ui{Bands below/above}] How many occupied and empty states around the
    gap receive a correction.
  \item[\ui{MPI ranks}] Yambo only.
\end{description}

\subsection{The viewer}

The \ui{GW Viewer} --- also on \menu{Results}\menusep\menu{GW Viewer\dots} ---
reports the DFT and quasiparticle band edges, the gap renormalization as its
headline number, and a per-state table of $\varepsilon_{\text{DFT}}$,
$E_{\text{qp}}$ and the correction between them.

\begin{calnote}
$G_0W_0$ essentially always \emph{opens} a semiconductor gap, typically by
0.5--2\,eV. A renormalization that is negative or near zero is far more likely
an unconverged screening cutoff or too few empty bands than a physical result,
and the viewer says so rather than presenting the number neutrally. The
per-state corrections are the other thing to read: they should vary smoothly
across bands, not jump erratically between neighbouring states.
\end{calnote}
